% =========================================================
%  config/listings.tex
%  Konfigurasi penampilan kode program: Python dan C# (Unity).
%  Seluruh buku memakai konfigurasi yang sama.
% =========================================================

% --- Gaya dasar kode program ----------------------------
\lstdefinestyle{bookstyle}{
  basicstyle=\ttfamily\footnotesize,
  breaklines=true,
  breakatwhitespace=false,
  postbreak=\mbox{\textcolor{ubAccent}{$\hookrightarrow$}\space},
  frame=single,
  framerule=0.4pt,
  rulecolor=\color{ubLine},
  numbers=left,
  numberstyle=\tiny\color{ubAccent},
  numbersep=9pt,
  showstringspaces=false,
  tabsize=4,
  keepspaces=true,
  columns=fullflexible,
  captionpos=b,
  xleftmargin=1.9em,
  framexleftmargin=1.5em,
  keywordstyle=\color{ubPrimary}\bfseries,
  commentstyle=\color{ubGreen}\itshape,
  stringstyle=\color{ubRed},
  emphstyle=\color{ubWarn}\bfseries,
  extendedchars=true,
  inputencoding=utf8
}

% --- Gaya keluaran program (tanpa nomor baris) ----------
\lstdefinestyle{outputstyle}{
  style=bookstyle,
  numbers=none,
  xleftmargin=1.1em,
  framexleftmargin=0.8em,
  basicstyle=\ttfamily\footnotesize\color{black},
  backgroundcolor=\color{ubGray}
}

% --- Bahasa Python --------------------------------------
\lstdefinelanguage{PythonBook}{
  language=Python,
  morekeywords={self,True,False,None,as,with,assert,lambda,yield,global,nonlocal},
  morecomment=[l]{\#},
  morestring=[b]"
}

% --- Bahasa C# (skrip Unity) ----------------------------
\lstdefinelanguage{CSharp}{
  morekeywords={
    abstract,as,base,bool,break,byte,case,catch,char,checked,class,const,
    continue,decimal,default,delegate,do,double,else,enum,event,explicit,
    extern,false,finally,fixed,float,for,foreach,goto,if,implicit,in,int,
    interface,internal,is,lock,long,namespace,new,null,object,operator,out,
    override,params,private,protected,public,readonly,ref,return,sbyte,
    sealed,short,sizeof,stackalloc,static,string,struct,switch,this,throw,
    true,try,typeof,uint,ulong,unchecked,unsafe,ushort,using,virtual,void,
    volatile,while,var,get,set,value,async,await
  },
  sensitive=true,
  morecomment=[l]{//},
  morecomment=[s]{/*}{*/},
  morestring=[b]",
  morestring=[b]'
}

% --- Istilah khusus Unity ditandai tersendiri -----------
\newcommand{\unityapi}[1]{\textcolor{ubWarn}{\textbf{#1}}}

\lstset{style=bookstyle}

% --- Perintah ringkas pemanggil kode --------------------
% \kodepython{python/ch04_game_loop.py}{Judul}{lst:label}
\newcommand{\kodepython}[3]{%
  \lstinputlisting[style=bookstyle, language=PythonBook,
    caption={#2}, label={#3}]{\codefile{#1}}%
}
% \kodecsharp{unity/ch13_player_controller.cs}{Judul}{lst:label}
\newcommand{\kodecsharp}[3]{%
  \lstinputlisting[style=bookstyle, language=CSharp,
    caption={#2}, label={#3}]{\codefile{#1}}%
}
% \kodekeluaran{...} untuk menampilkan potongan keluaran program
\lstnewenvironment{kodekeluaran}[1][]{%
  \lstset{style=outputstyle, #1}}{}
